\documentclass[12pt,a4paper]{article}

% ============================================================
%  PACKAGES
% ============================================================
\usepackage[utf8]{inputenc}
\usepackage[T1]{fontenc}
\usepackage[french]{babel}
\usepackage{amsmath, amssymb}
\usepackage{geometry}
\usepackage{booktabs}
\usepackage{array}
\usepackage{xcolor}
\usepackage{hyperref}
\usepackage{graphicx}
\usepackage{fancyhdr}
\usepackage{titlesec}
\usepackage{enumitem}
\usepackage{listings}
\usepackage{float}

\geometry{top=2.5cm, bottom=2.5cm, left=3cm, right=3cm}

% ============================================================
%  COULEURS
% ============================================================
\definecolor{bleuLIU}{RGB}{0, 51, 102}
\definecolor{grisLeger}{RGB}{245, 245, 245}
\definecolor{vertCode}{RGB}{0, 128, 0}
\definecolor{rougeCode}{RGB}{180, 0, 0}

% ============================================================
%  STYLE DES SECTIONS
% ============================================================
\titleformat{\section}
  {\large\bfseries\color{bleuLIU}}
  {\thesection.}{0.5em}{}[\titlerule]

\titleformat{\subsection}
  {\normalsize\bfseries\color{bleuLIU}}
  {\thesubsection.}{0.5em}{}

% ============================================================
%  EN-TÊTE ET PIED DE PAGE
% ============================================================
\pagestyle{fancy}
\fancyhf{}
\rhead{\textcolor{bleuLIU}{\small Recherche Opérationnelle — 2025/2026}}
\lhead{\textcolor{bleuLIU}{\small Méthode du Simplexe}}
\cfoot{\thepage}
\renewcommand{\headrulewidth}{0.4pt}

% ============================================================
%  STYLE DU CODE PYTHON
% ============================================================
\lstset{
  language=Python,
  backgroundcolor=\color{grisLeger},
  basicstyle=\ttfamily\small,
  keywordstyle=\color{bleuLIU}\bfseries,
  commentstyle=\color{vertCode}\itshape,
  stringstyle=\color{rougeCode},
  frame=single,
  rulecolor=\color{bleuLIU},
  numbers=left,
  numberstyle=\tiny\color{gray},
  breaklines=true,
  showstringspaces=false,
  tabsize=4
}

% ============================================================
%  PAGE DE TITRE
% ============================================================
\begin{document}

\begin{titlepage}
  \centering
  \vspace*{1cm}

  {\Large \textbf{Université Libanaise Internationale en Mauritanie}}\\[0.3cm]
  {\large Faculté des Sciences}\\[0.1cm]
  {\large Module : Recherche Opérationnelle — Licence}\\[0.1cm]
  {\large Responsable : Dr. EL BENANY Med Mahmoud}

  \vspace{1.5cm}
  \rule{\textwidth}{1.5pt}\\[0.4cm]
  {\LARGE \textbf{\textcolor{bleuLIU}{Rapport de Projet}}}\\[0.3cm]
  {\huge \textbf{Méthode du Simplexe}}\\[0.3cm]
  {\large Sujet 1.4 — Sujets classiques et fondamentaux}\\[0.2cm]
  \rule{\textwidth}{1.5pt}

  \vspace{1.5cm}

  \begin{tabular}{rl}
    \textbf{Étudiant :}   & \textit{[Votre Nom Complet]} \\[0.3cm]
    \textbf{Filière :}    & Licence Informatique / Mathématiques \\[0.3cm]
    \textbf{Année :}      & 3\textsuperscript{ème} année \\[0.3cm]
    \textbf{Année univ. :}& 2025 -- 2026 \\
  \end{tabular}

  \vfill
  {\large Nouakchott, Mai 2026}
\end{titlepage}

% ============================================================
%  TABLE DES MATIÈRES
% ============================================================
\tableofcontents
\newpage

% ============================================================
%  1. INTRODUCTION
% ============================================================
\section{Introduction}

La \textbf{programmation linéaire} (PL) est l'une des méthodes les plus 
fondamentales de la recherche opérationnelle. Elle permet de modéliser 
des problèmes de décision réels (allocation de ressources, planification, 
logistique) sous forme d'un programme mathématique avec une fonction 
objectif linéaire et des contraintes linéaires.

La \textbf{méthode du simplexe}, développée par George Dantzig en 1947, 
est l'algorithme de référence pour résoudre ces problèmes. Elle explore 
les sommets du polyèdre des solutions réalisables de façon itérative 
jusqu'à atteindre l'optimum global.

\medskip
Ce rapport présente :
\begin{itemize}[noitemsep]
  \item la formulation mathématique du problème,
  \item la mise en forme standard,
  \item la résolution manuelle étape par étape,
  \item l'implémentation en Python avec la bibliothèque \texttt{PuLP},
  \item l'analyse de la solution optimale et la notion de dualité.
\end{itemize}

% ============================================================
%  2. FORMULATION MATHÉMATIQUE
% ============================================================
\section{Formulation mathématique}

\subsection{Problème original}

Nous considérons le problème de programmation linéaire suivant :

\[
  \max \; z = 3x + 4y
\]
\[
  \text{sous les contraintes :} \quad
  \begin{cases}
    x + 2y \leq 8 \\
    3x + y  \leq 9 \\
    x, y \geq 0
  \end{cases}
\]

\textbf{Interprétation économique :} On cherche à maximiser un profit $z$ 
en choisissant les quantités $x$ et $y$ de deux produits, sous des 
contraintes de disponibilité de ressources (heures-machines, matières 
premières, etc.).

\subsection{Mise en forme standard}

Pour appliquer le simplexe, on transforme les inégalités en égalités en 
ajoutant des \textbf{variables d'écart} $s_1, s_2 \geq 0$ :

\[
  \max \; z = 3x + 4y + 0s_1 + 0s_2
\]
\[
  \begin{cases}
    x + 2y + s_1 \phantom{+ s_2} = 8 \\
    3x + y \phantom{+ s_1} + s_2 = 9 \\
    x, y, s_1, s_2 \geq 0
  \end{cases}
\]

La \textbf{solution de base initiale} (SBI) est obtenue en posant $x = y = 0$ :
\[
  s_1 = 8, \quad s_2 = 9, \quad z = 0
\]

% ============================================================
%  3. RÉSOLUTION MANUELLE
% ============================================================
\section{Résolution manuelle par le simplexe}

\subsection{Tableau initial}

Le tableau simplex initial est construit à partir de la forme standard. 
La ligne $z$ représente $z - 3x - 4y = 0$, donc les coefficients de la 
fonction objectif apparaissent avec un signe négatif :

\begin{table}[H]
\centering
\caption{Tableau simplex — Itération 0 (initial)}
\begin{tabular}{c|cccc|c}
\toprule
\textbf{VB} & $x$ & $y$ & $s_1$ & $s_2$ & \textbf{RHS} \\
\midrule
$s_1$ & 1 & \textbf{2} & 1 & 0 & 8 \\
$s_2$ & 3 & 1 & 0 & 1 & 9 \\
\midrule
$z$   & $-3$ & $-4$ & 0 & 0 & 0 \\
\bottomrule
\end{tabular}
\end{table}

\textbf{Choix de la variable entrante :} Le coefficient le plus négatif 
dans la ligne $z$ est $-4$ (colonne $y$) $\Rightarrow$ \textbf{variable 
entrante : $y$} (colonne pivot).

\textbf{Choix de la variable sortante (test du rapport minimum) :}
\[
  \min\left\{\frac{8}{2},\; \frac{9}{1}\right\} = \min\{4,\; 9\} = 4 
  \Rightarrow \text{ligne } s_1 \Rightarrow \textbf{variable sortante : } s_1
\]
\textbf{Élément pivot : } $a_{11} = 2$ (ligne $s_1$, colonne $y$).

\subsection{Itération 1}

On divise la ligne pivot ($s_1$) par l'élément pivot (2), puis on 
élimine $y$ des autres lignes par combinaisons linéaires :

\begin{table}[H]
\centering
\caption{Tableau simplex — Itération 1}
\begin{tabular}{c|cccc|c}
\toprule
\textbf{VB} & $x$ & $y$ & $s_1$ & $s_2$ & \textbf{RHS} \\
\midrule
$y$   & $\frac{1}{2}$ & 1 & $\frac{1}{2}$ & 0 & 4 \\
$s_2$ & \textbf{$\frac{5}{2}$} & 0 & $-\frac{1}{2}$ & 1 & 5 \\
\midrule
$z$   & $-1$ & 0 & $2$ & 0 & 16 \\
\bottomrule
\end{tabular}
\end{table}

\textbf{Variable entrante :} coefficient $-1$ (colonne $x$) $\Rightarrow$ 
\textbf{variable entrante : $x$}.

\textbf{Test du rapport minimum :}
\[
  \min\left\{\frac{4}{1/2},\; \frac{5}{5/2}\right\} = \min\{8,\; 2\} = 2
  \Rightarrow \text{ligne } s_2 \Rightarrow \textbf{variable sortante : } s_2
\]
\textbf{Élément pivot : } $\frac{5}{2}$ (ligne $s_2$, colonne $x$).

\subsection{Itération 2 — Solution optimale}

On divise la ligne pivot par $\frac{5}{2}$, puis on élimine $x$ :

\begin{table}[H]
\centering
\caption{Tableau simplex — Itération 2 (optimal)}
\begin{tabular}{c|cccc|c}
\toprule
\textbf{VB} & $x$ & $y$ & $s_1$ & $s_2$ & \textbf{RHS} \\
\midrule
$y$ & 0 & 1 & $\frac{2}{5}$ & $-\frac{1}{5}$ & 3 \\
$x$ & 1 & 0 & $-\frac{1}{5}$ & $\frac{2}{5}$ & 2 \\
\midrule
$z$ & 0 & 0 & $\frac{9}{5}$ & $\frac{2}{5}$ & \textbf{18} \\
\bottomrule
\end{tabular}
\end{table}

Tous les coefficients de la ligne $z$ sont $\geq 0$ : 
\textbf{la solution est optimale}.

\[
  \boxed{x^* = 2, \quad y^* = 3, \quad z^* = 18}
\]

\textbf{Vérification :}
\[
  x + 2y = 2 + 6 = 8 \leq 8 \;\checkmark \qquad 
  3x + y = 6 + 3 = 9 \leq 9 \;\checkmark
\]

% ============================================================
%  4. IMPLÉMENTATION PYTHON
% ============================================================
\section{Implémentation Python avec PuLP}

\subsection{Code source}

\begin{lstlisting}[caption={Résolution du problème par PuLP}]
from pulp import *

# Création du problème
prob = LpProblem("Simplexe_1_4", LpMaximize)

# Variables de décision
x = LpVariable("x", lowBound=0)
y = LpVariable("y", lowBound=0)

# Fonction objectif
prob += 3*x + 4*y, "Profit"

# Contraintes
prob += x + 2*y <= 8, "Contrainte_1"
prob += 3*x + y  <= 9, "Contrainte_2"

# Résolution
prob.solve(PULP_CBC_CMD(msg=0))

# Affichage des résultats
print("=== SOLUTION OPTIMALE ===")
print(f"Statut     : {LpStatus[prob.status]}")
print(f"x*         : {value(x)}")
print(f"y*         : {value(y)}")
print(f"z* (profit): {value(prob.objective)}")

print("\n=== PRIX FANTOMES (Dualite) ===")
for name, c in prob.constraints.items():
    print(f"{name} : shadow price = {c.pi:.4f}")
\end{lstlisting}

\subsection{Résultats}

\begin{table}[H]
\centering
\caption{Résultats de la résolution Python}
\begin{tabular}{lc}
\toprule
\textbf{Variable / Indicateur} & \textbf{Valeur} \\
\midrule
Statut & Optimal \\
$x^*$ & 2.0 \\
$y^*$ & 3.0 \\
$z^*$ & 18.0 \\
Prix fantôme — Contrainte 1 & 1.8 \\
Prix fantôme — Contrainte 2 & 0.4 \\
\bottomrule
\end{tabular}
\end{table}

Les résultats Python confirment exactement la résolution manuelle.

% ============================================================
%  5. ANALYSE ET DUALITÉ
% ============================================================
\section{Analyse de la solution et dualité}

\subsection{Interprétation de la solution optimale}

La solution optimale $x^* = 2$, $y^* = 3$ signifie qu'il faut 
produire \textbf{2 unités} du produit $x$ et \textbf{3 unités} du 
produit $y$ pour maximiser le profit à \textbf{18 unités monétaires}.

On remarque que :
\begin{itemize}[noitemsep]
  \item La contrainte 1 est \textbf{saturée} ($x + 2y = 8$) : 
  la ressource 1 est entièrement utilisée.
  \item La contrainte 2 est \textbf{saturée} ($3x + y = 9$) : 
  la ressource 2 est entièrement utilisée.
\end{itemize}

Les deux variables d'écart valent $s_1 = s_2 = 0$ : il n'y a 
\textbf{aucune ressource inutilisée} à l'optimum.

\subsection{Prix fantômes (Variables duales)}

Les \textbf{prix fantômes} (shadow prices) mesurent l'amélioration 
marginale du profit optimal si l'on augmente d'une unité le 
membre droit d'une contrainte.

\begin{table}[H]
\centering
\caption{Prix fantômes et interprétation}
\begin{tabular}{clcc}
\toprule
\textbf{Contrainte} & \textbf{Description} & \textbf{Prix fantôme} & \textbf{Impact} \\
\midrule
$x + 2y \leq 8$ & Ressource 1 & $y_1^* = \dfrac{9}{5} = 1{,}8$ 
  & $+1$ unité $\Rightarrow z^* = 19{,}8$ \\[8pt]
$3x + y \leq 9$  & Ressource 2 & $y_2^* = \dfrac{2}{5} = 0{,}4$ 
  & $+1$ unité $\Rightarrow z^* = 18{,}4$ \\
\bottomrule
\end{tabular}
\end{table}

\textbf{Conclusion managériale :} Si on devait investir pour augmenter 
une seule ressource, il vaut mieux augmenter la \textbf{ressource 1} 
(prix fantôme $= 1{,}8 > 0{,}4$) car elle génère un gain plus important.

\subsection{Problème dual}

Le problème dual associé est :

\[
  \min \; w = 8y_1 + 9y_2
\]
\[
  \begin{cases}
    y_1 + 3y_2 \geq 3 \\
    2y_1 + y_2 \geq 4 \\
    y_1, y_2 \geq 0
  \end{cases}
\]

Par le \textbf{théorème des écarts complémentaires}, la solution duale 
optimale est $y_1^* = \frac{9}{5}$ et $y_2^* = \frac{2}{5}$, 
avec $w^* = 8 \times \frac{9}{5} + 9 \times \frac{2}{5} = 18 = z^*$.

Ce résultat confirme le \textbf{théorème de dualité forte} : 
les valeurs optimales des problèmes primal et dual sont égales.

% ============================================================
%  6. CONCLUSION
% ============================================================
\section{Conclusion}

Ce projet a permis de maîtriser la méthode du simplexe dans ses deux 
dimensions : théorique et pratique.

Sur le plan \textbf{théorique}, nous avons vu comment formuler un 
problème de PL en forme standard, appliquer les itérations du simplexe 
et interpréter la solution à travers la dualité et les prix fantômes.

Sur le plan \textbf{pratique}, l'implémentation Python avec 
\texttt{PuLP} confirme les résultats manuels et ouvre la voie à la 
résolution de problèmes de plus grande dimension inaccessibles 
manuellement.

La méthode du simplexe constitue un socle essentiel pour aborder 
des méthodes plus avancées comme le \textit{branch-and-bound} 
(pour les variables entières) ou les algorithmes de points intérieurs.

\end{document}
